\chapter{Competition and the Future}\label{ch:fm:competition-and-the-future}

When the US securities regulator proposed in December 2022 to make many retail orders compete in auctions, its own release found that broker-dealers
sent more than 90\% of individual investors' marketable orders in listed stocks to six off-exchange dealers, the wholesalers, and that two of them
handled about 66\% of the wholesalers' volume. The proposal was withdrawn in June 2025. Trading looks like a market of many firms; business by
business it is often a market of a few. This last chapter asks why, what protects the few, and what three changes now under way (tokenised assets,
round-the-clock markets and machine-learning tooling) do to the answer.

\section{Moats in trading}

\begin{definition}[Economic moat]\label{def:fm:competition-and-the-future:moat}
An \emph{economic moat}\index{economic moat} is a lasting advantage that lets a firm earn returns above its cost of capital in a business without
competitors entering until those returns are gone: a cost advantage from scale, proprietary information or order flow, customers' \omterm{def:fm:build-against-buy:lockin}{switching costs},
network effects, or a licence or relationship others cannot obtain quickly.
\end{definition}

Book 11 showed that a trading edge decays as others find it: capture falls as competitors enter. The durable advantages are therefore rarely the
signals. They are the fixed costs a new entrant must pay before its first trade (chapter 1's \omterm{def:fm:the-economics-of-a-trading-firm:fixed}{fixed cost base}, chapter 2's \omterm{def:fm:the-proprietary-market-making-firm:treadmill}{technology treadmill},
chapter 19's \omterm{def:fm:technology-strategy:race}{latency arms race}), the order flow that comes through relationships and contracts (chapter 23), and the data and history a firm has
accumulated. Sutton's work on sunk costs and market structure separates two cases. When the fixed cost of entry is exogenous, a larger market
supports more firms and concentration falls as it grows. When firms can spend to be better (faster, better data, better models) and customers
reward the best, the spending itself rises with the market, and concentration need not fall at all.

\begin{definition}[Market concentration, Herfindahl--Hirschman index]\label{def:fm:competition-and-the-future:hhi}
\emph{Market concentration}\index{market concentration} is the degree to which a business's volume or revenue is held by a few firms, measured by
concentration ratios (the $k$ largest firms' combined share, $\mathrm{CR}_k$) or by the \emph{Herfindahl--Hirschman index}\index{Herfindahl--Hirschman
index}, the sum of the squared market shares in percent,
\[ \mathrm{HHI}=\sum_{i=1}^n(100\,s_i)^2, \]
from near 0 for many small firms to 10\,000 for one (Book 1's Herfindahl index of venues, on a 0 to 10\,000 scale).
\end{definition}

\section{Consolidation and market structure}

The public record rarely gives every firm's share; the release gives two numbers. They bound the index.

\begin{proposition}[Bounds from partial shares]\label{prop:fm:competition-and-the-future:bounds}
If a business has $n$ firms whose $k$ largest hold a combined share $S_k$, the index is smallest when the $k$ leaders are equal and the other firms
are equal, $\mathrm{HHI}_{\min}=10^4\bigl(S_k^2/k+(1-S_k)^2/(n-k)\bigr)$, provided $(1-S_k)/(n-k)\le S_k/k$. For $n=6$, $k=2$ and $S_2=0.66$ it
lies between 2\,467 and 3\,668.
\end{proposition}

\begin{proof}
For a fixed total, a sum of squares is smallest when the parts are equal; applying this to the leaders and to the others separately gives the
minimum, which is feasible when each other firm is no larger than each leader. For the maximum, fix the smaller leader's share $b$: the larger
leader holds $S_2-b$, and the others, each at most $b$, hold $1-S_2$; with $b=(1-S_2)/4=0.085$ all four others must equal $b$, and the leader
holds 0.575, giving $10^4(0.575^2+5\times0.085^2)=3\,668$. A search over $b$ (\texttt{firm.moats.hhi\_bounds}) confirms that larger $b$ gives less.
\end{proof}

\begin{omfigure}
\begin{tikzpicture}
\begin{axis}[width=10.5cm,height=5cm,ybar,bar width=10pt,ymin=0,ymax=65,ylabel={\small share (\%)},xtick={1,2,3,4,5,6},
  xlabel={\small wholesaler, by rank},tick label style={font=\footnotesize},grid=major,grid style={gray!20},table/col sep=comma,
  legend style={font=\scriptsize,at={(0.97,0.95)},anchor=north east},legend cell align=left]
\addplot[fill=omDef!60,draw=omDef] table[x=rank,y=low]{figdata/desk/30-competition-and-the-future/bounds.csv};
\addplot[fill=omThm!45,draw=omThm] table[x=rank,y=high]{figdata/desk/30-competition-and-the-future/bounds.csv};
\legend{least concentrated (index 2\,467),most concentrated (index 3\,668)}
\end{axis}
\end{tikzpicture}
\omcaption{The two share vectors that bound the index for six wholesalers whose two largest hold 66\%, the figures in the SEC's 2022 release. Either
way the business is highly concentrated under the 2023 merger guidelines' threshold of 1\,800. Data: \texttt{fm\_competition.public\_bounds}.}\label{fig:fm:competition-and-the-future:bounds}
\end{omfigure}

\begin{dated}{2026-09}{The public record of the chapter}\label{dat:fm:competition-and-the-future:record}
\textbf{Wholesaling}: the SEC's Order Competition Rule proposal (Release 34-96495, published 3 January 2023) found more than 90\% of individual
investors' marketable orders in NMS stocks routed to six wholesalers, two with about 66\% of the wholesalers' executed share volume in the first
quarter of 2022; the Commission withdrew the proposal on 17 June 2025. \textbf{Merger guidelines} (DOJ and FTC, December 2023): an index above 1\,800 is
highly concentrated; an increase above 100 in such a market, or a merged share above 30\% with an increase above 100, is presumed to lessen
competition substantially. \textbf{\omterm{def:fm:competition-and-the-future:token}{Tokenised fund}}: a registered US government money-market fund's prospectus of 1 August 2026 states that its
transfer agent keeps the official record of share ownership in a blockchain-integrated system on public networks. \textbf{Round-the-clock}: the SEC
registered 24X National Exchange on 27 November 2024 to trade listed stocks 23 hours a day, five days a week; NSCC extended its clearing hours to
24x5 on 29 June 2026, with exchanges expected to follow in late 2026.
\end{dated}

\begin{definition}[Free-entry equilibrium]\label{def:fm:competition-and-the-future:freeentry}
A \emph{free-entry equilibrium}\index{free-entry equilibrium} is the number of firms at which each firm in a business covers its fixed cost and one
more entrant would not: $\pi(n)\ge F>\pi(n+1)$, where $\pi(n)$ is each firm's operating profit when $n$ firms compete and $F$ is the fixed cost of
being in the business.
\end{definition}

In a Cournot market with linear demand $a-bQ$, a common marginal cost $c$ and $n$ identical firms, each earns $\pi(n)=S/(n+1)^2$ with
$S=(a-c)^2/b$, so the free-entry number is $n^*=\lfloor\sqrt{S/F}\rfloor-1$ and the index is $10\,000/n^*$. Concentration is set by the ratio of the
market's size to the fixed cost, not by either alone.

\section{Tutorial: the fixed cost falls}\label{tut:fm:competition-and-the-future:tut}

\begin{tutorial}
\textbf{Goal.} Compute the free-entry number of firms and the index at the book's parameters; halve and double the fixed cost; merge two firms, with
and without synergies; and bound the index from the public figures. \textbf{End state:}
\cref{fig:fm:competition-and-the-future:entry,fig:fm:competition-and-the-future:merger}.
\begin{tutsteps}
\item \textbf{Parameters.} $S=\$3\,000$ million a year (demand intercept 10 and marginal cost 4, in basis points), and a fixed cost $F=\$50$ million
a year per firm; both are illustrative.
\item \textbf{Free entry.} \texttt{free\_entry(3000, 50)} gives six firms, each earning \$61.2 million before and \$11.2 million after its fixed cost;
the index is 1\,667.
\item \textbf{Halve and double.} At \$25 million, nine firms (1\,111); at \$100 million, four (2\,500).
\item \textbf{The planner.} \texttt{planner} maximises consumers' surplus plus profits less fixed costs: three firms at \$50 million.
\item \textbf{A merger.} Two of the six merge and save one fixed cost; then find the marginal-cost synergy that keeps the price unchanged.
\end{tutsteps}
\end{tutorial}

\omcode{code/firm/moats/firm_moats.py}{33}{64}{Bounds on the index from partial shares, the free-entry number of firms, total surplus and the planner's
number.}\label{lst:fm:competition-and-the-future:entry}

\begin{omfigure}
\begin{tikzpicture}
\begin{axis}[width=10.5cm,height=5.6cm,xlabel={\small fixed cost per firm (\$ million a year)},ylabel={\small index (HHI)},xmin=10,xmax=200,ymin=0,
  ymax=6000,tick label style={font=\footnotesize},grid=major,grid style={gray!20},table/col sep=comma,legend cell align=left,
  legend style={font=\scriptsize,at={(0.03,0.97)},anchor=north west},scaled y ticks=false,
  yticklabel style={/pgf/number format/.cd,1000 sep={,}}]
\fill[omThm!12] (axis cs:10,2467) rectangle (axis cs:200,3668);
\addplot[omDef,very thick,const plot] table[x=F,y=hhi]{figdata/desk/30-competition-and-the-future/entry.csv};
\draw[gray,dashed] (axis cs:10,1800) -- (axis cs:200,1800) node[pos=0.8,below,font=\scriptsize] {1\,800: highly concentrated};
\node[font=\scriptsize,omThm,anchor=west] at (axis cs:14,3070) {wholesalers, 2022: bounds};
\draw[black] (axis cs:50,1666.7) circle[radius=2.5pt];
\legend{free-entry index}
\end{axis}
\end{tikzpicture}
\omcaption{The free-entry index against the fixed cost at $S=\$3\,000$ million: six firms and 1\,667 at \$50 million (circled), nine and 1\,111
at \$25 million, four and 2\,500 at \$100 million. The shaded band is the range the public wholesaler figures allow. Data:
\texttt{fm\_competition.entry}.}\label{fig:fm:competition-and-the-future:entry}
\end{omfigure}

The symmetric model at the book's parameters gives six firms and an index of 1\,667, below the guidelines' threshold; the public wholesaler figures
imply at least 2\,467. The gap is the model's symmetry: the real firms differ in cost, and a firm with lower costs takes a larger share. Halving the
fixed cost, as cloud computing and open-source tooling do for research and infrastructure, raises the free-entry number from six to nine; doubling it,
as a new speed tier does for the firms that must buy it, lowers it to four.

\begin{remark}[Excess entry]\label{rem:fm:competition-and-the-future:excess}
The planner would choose three firms at \$50 million, not six: each entrant takes business from the others, and its private gain exceeds the
surplus it adds. Total surplus is \$1\,169 million at free entry against \$1\,256 million at three firms. This is Mankiw and Whinston's excess-entry
result in the chapter's model; it is not an argument for fewer firms in any real business, where entrants also bring new products and lower costs.
\end{remark}

Consolidation in trading has followed the arithmetic of fixed costs. Knight Capital and GETCO merged under KCG Holdings on 1 July 2013; Virtu
Financial acquired KCG on 20 July 2017 for \$20.00 a share in cash, and Investment Technology Group on 1 March 2019 for about \$1.0 billion
(\cref{fig:fm:competition-and-the-future:chain}). Each deal combined firms whose fixed costs (technology, connectivity, data, compliance)
overlapped; chapter 20 read the second as a build-against-buy decision.

\begin{omfigure}
\begin{tikzpicture}[font=\small,firm/.style={draw=omDef,fill=omDef!8,rounded corners,align=center,minimum height=0.8cm,text width=2.6cm},
  arr/.style={-{Stealth[length=2mm]},thick,omDef}]
\node[firm] (kn) at (0,1.4) {Knight Capital};
\node[firm] (ge) at (0,0) {GETCO};
\node[firm] (kcg) at (4,0.7) {KCG Holdings\\{\scriptsize 1 July 2013}};
\node[firm] (vi) at (9.2,0.7) {Virtu Financial\\{\scriptsize KCG: 20 July 2017}};
\node[firm] (itg) at (9.2,-1.2) {Investment\\Technology Group};
\draw[arr] (kn) -- (kcg);
\draw[arr] (ge) -- (kcg);
\draw[arr] (kcg) -- node[above,font=\scriptsize] {\$20.00 a share} (vi);
\draw[arr] (itg) -- node[right,font=\scriptsize,align=left] {1 March 2019\\about \$1.0 billion} (vi);
\end{tikzpicture}
\omcaption{A consolidation chain from the public filings: two firms merged in 2013, the combined firm was acquired in 2017, and the acquirer bought
Investment Technology Group in 2019.}\label{fig:fm:competition-and-the-future:chain}
\end{omfigure}

\begin{omfigure}
\begin{tikzpicture}
\begin{axis}[width=10.5cm,height=5.2cm,xlabel={\small marginal-cost synergy of the merged firm (bp)},ylabel={\small price (bp)},xmin=0,xmax=1.2,
  ymin=4.78,ymax=5.02,tick label style={font=\footnotesize},grid=major,grid style={gray!20},table/col sep=comma,
  yticklabel style={/pgf/number format/.cd,fixed,precision=2}]
\addplot[omDef,very thick] table[x=synergy,y=price]{figdata/desk/30-competition-and-the-future/merger.csv};
\draw[omThm,dashed] (axis cs:0,4.857) -- (axis cs:1.2,4.857) node[pos=0.05,below right,font=\scriptsize] {price before the merger};
\draw[gray,dotted] (axis cs:0.857,4.78) -- (axis cs:0.857,5.02) node[pos=0.93,right,font=\scriptsize] {0.857};
\end{axis}
\end{tikzpicture}
\omcaption{Two of six Cournot firms merge: the price after the merger against the merged firm's reduction in marginal cost. Without it, the price
rises from 4.857 to 5.000 (2.9\%); a reduction of 0.857, 21.4\% of the marginal cost, restores the old price. Data:
\texttt{fm\_competition.merger}.}\label{fig:fm:competition-and-the-future:merger}
\end{omfigure}

\omcode{code/firm/moats/firm_moats.py}{67}{90}{Cournot equilibrium with unequal costs, a merger with a marginal-cost synergy, and the synergy that keeps
the price unchanged.}\label{lst:fm:competition-and-the-future:merge}

The merger in the model shows why consolidation happens even when it raises prices. Without a marginal-cost synergy, the price rises 2.9\% and the
index from 1\,667 to 2\,000, but the merged firm saves one fixed cost: the two firms' profit after fixed costs rises from \$22.4 million to \$33.3
million, while consumers lose \$60 million of surplus. With the synergy that holds the price, the merged firm's share is 33.3\% and the index 2\,222,
above both of the guidelines' presumptions, although buyers pay no more: concentration measures structure, not harm, which is why the guidelines
treat their thresholds as presumptions to be tested.

\section{Tokenised assets}

\begin{definition}[Asset tokenisation, tokenised fund]\label{def:fm:competition-and-the-future:token}
\emph{Asset tokenisation}\index{asset tokenisation} records ownership of a financial asset (a fund share, a bond, a deposit) as tokens on a
blockchain, with the issuer, its transfer agent or a \omterm{def:fm:the-asset-manager-and-the-fund:admin}{custodian} keeping legal responsibility for the record. A \emph{tokenised fund}\index{tokenised
fund} is a fund whose official register of share ownership is kept, wholly or partly, in such a system.
\end{definition}

The registered money-market fund in the dated box shows what the record currently looks like: its transfer agent keeps the official register in a
permissioned system on public blockchains, with a whitelist of approved wallets in a smart contract (Book 3), the power to correct errors and to limit
transfers, and peer-to-peer transfers of shares between approved wallets. For a trading firm the change is operational before it is strategic. A
share that can move between wallets at any hour can serve as collateral outside banking hours, which bears on chapter 14's margin and funding; its
records sit with a transfer agent that can reverse transfers, so it is not a bearer asset. Whether tokenisation lowers the fixed cost of entering a
business (by removing intermediaries) or raises it (new custody, controls and connections) is the question the model above makes precise.

\section{Round-the-clock markets}

\begin{definition}[Round-the-clock trading]\label{def:fm:competition-and-the-future:rtc}
\emph{Round-the-clock trading}\index{round-the-clock trading} extends a market's trading sessions to most or all hours of the week, in the listed
equity case to 23 or 24 hours on business days, with the clearing, market data and surveillance that the sessions need running over the same hours.
\end{definition}

The Commission's order registering 24X National Exchange describes the structure: the usual pre-market, core and post-market sessions from 4 a.m. to
7 p.m., and an overnight session from 8 p.m. to 4 a.m. on the nights before business days, weekends and holidays having been removed from the
proposal. The clearing house moved first: NSCC's clearing runs from Sunday 8 p.m. to Friday 8 p.m. For a firm, longer hours raise fixed costs
(staffing, monitoring, incident response through the night, chapter 21's on-call design) while the overnight volume, at first, is thin. In the
chapter's model that is a rise in $F$ for every firm that chooses to be present, and a smaller $S$ per session: the overnight business will support
fewer firms than the day.

\section{What machine-learning tooling changes}

Book 12's machine-learning platforms and large language models cut two costs at once. The cost of research infrastructure falls: open-source libraries,
rented accelerators and pretrained models do what a team once built. So does the cost of routine work (reconciliations, documentation, code review,
first drafts of research), chapter 21's organisation with fewer people per function. Both lower $F$, and the free-entry model says more firms. Sutton's
second case says to look further: if the tools also let the best firms spend to be better (more data, larger models, faster iteration), and the
business rewards being best, the spending rises with the market and the leaders' share need not fall. The two effects can be told apart in a firm's
own numbers: which costs fell (chapter 1's \omterm{def:fm:the-economics-of-a-trading-firm:fixed}{fixed cost base}), and whether its capture held against the best competitor or against the average one.

\begin{method}[Reading a business's structure]\label{met:fm:competition-and-the-future:read}
\begin{enumerate}
\item Measure concentration from what is public: counts, top shares, bounds on the index.
\item Estimate the business's size and the fixed cost of being in it; compare the free-entry number with the firms observed.
\item Ask which costs are exogenous and which the firms choose to spend to be better; the second kind sustains concentration as the market grows.
\item For each change (a tool, a market structure, a rule), ask whether it moves $S$, $F$ or the spending race, and for whom.
\end{enumerate}
\end{method}

\section{Build: the moats model}\label{bld:fm:competition-and-the-future:moats}

\begin{build}
\bfield{Purpose} Concentration from shares or partial shares, the free-entry number of firms, the planner's number, and Cournot mergers with fixed-
and marginal-cost synergies.
\bfield{Interface} \texttt{firm.moats}: \texttt{hhi}, \texttt{cr}, \texttt{hhi\_bounds}, \texttt{free\_entry}, \texttt{welfare}, \texttt{planner},
\texttt{cournot}, \texttt{merge}, \texttt{synergy\_for\_price}, \texttt{scenarios}.
\bfield{Rules} Shares are normalised; a Cournot firm with no output raises an error rather than silently dropping out; the bounds reject share data
that cannot be ordered.
\bfield{Acceptance tests} \texttt{code/firm/moats/tests/}: the index on hand examples; bounds that bracket random feasible share vectors; the
free-entry condition and excess entry; the merger's price rise and the synergy that removes it.
\bfield{Stretch} Endogenous sunk costs (firms choose quality at a cost); entry with asymmetric costs; two sessions (day and overnight) with different
sizes and a shared fixed cost.
\end{build}

\begin{omsources}
\item US Securities and Exchange Commission, Order Competition Rule, Release 34-96495, 88 FR 128 (2023); Withdrawal of Proposed Regulatory Actions,
90 FR 25531 (2025).
\item US Department of Justice and Federal Trade Commission, \emph{Merger Guidelines}, 18 December 2023.
\item KCG Holdings, Form 8-K12G3, 1 July 2013; Virtu Financial, Forms 8-K, 21 July 2017 and 1 March 2019.
\item Franklin OnChain U.S. Government Money Fund, prospectus, 1 August 2026.
\item SEC, In the Matter of the Application of 24X National Exchange LLC, Release 34-101777, 27 November 2024; DTCC, press release, 29 June 2026.
\item J. Sutton, \emph{Sunk Costs and Market Structure}, MIT Press, 1991.
\item N.~G. Mankiw and M.~D. Whinston, ``Free entry and social inefficiency'', \emph{RAND Journal of Economics}, 1986.
\end{omsources}

\section{Exercises}

\begin{exercise}[$\star$]\label{exo:fm:competition-and-the-future:1}
Compute the index of six equal firms, and of shares 33, 33, 8.5, 8.5, 8.5 and 8.5\%.
\end{exercise}

\begin{exercise}[$\star$]\label{exo:fm:competition-and-the-future:2}
Define an \omterm{def:fm:competition-and-the-future:moat}{economic moat} and give three sources of one in trading.
\end{exercise}

\begin{exercise}[$\star$]\label{exo:fm:competition-and-the-future:3}
With $S=\$3\,000$ million, how many firms does free entry support at a fixed cost of \$25 million, \$50 million and \$100 million?
\end{exercise}

\begin{exercise}[$\star\star$]\label{exo:fm:competition-and-the-future:4}
Prove \cref{prop:fm:competition-and-the-future:bounds}. Why must each other firm be no larger than each leader?
\end{exercise}

\begin{exercise}[$\star\star$]\label{exo:fm:competition-and-the-future:5}
Two of the six firms merge without a marginal-cost synergy. Would the merger be presumed to lessen competition under the 2023 guidelines, and is
it profitable for the merging firms?
\end{exercise}

\begin{exercise}[$\star\star$]\label{exo:fm:competition-and-the-future:6}
How does overnight trading change a market-making firm's fixed costs, and what does the free-entry model predict for the overnight session?
\end{exercise}

\begin{exercise}[$\star\star\star$]\label{exo:fm:competition-and-the-future:7}
\emph{Coding.} Compute total surplus at free entry and at the planner's number for $F=\$50$ million. By how much does free entry fall short?
\end{exercise}

\begin{exercise}[$\star\star\star$]\label{exo:fm:competition-and-the-future:8}
\emph{Find the flaw.} ``Machine-learning tools make research cheap, so the business will become competitive and the leaders' share will fall.''
\end{exercise}

\section{Problem: The Fixed Cost Falls}

\begin{problem}[{Weekend problem --- the fixed cost falls}]\label{pb:fm:competition-and-the-future:1}
A firm's strategy committee asks what cheaper research tools, tokenised collateral and round-the-clock markets will do to the businesses it is in.

\medskip\noindent\textbf{Part I --- The concepts.}
\begin{enumerate}
\item Define an \omterm{def:fm:competition-and-the-future:moat}{economic moat}.
\item Define \omterm{def:fm:competition-and-the-future:hhi}{market concentration} and the \omterm{def:fm:competition-and-the-future:hhi}{Herfindahl--Hirschman index}.
\item Define a \omterm{def:fm:competition-and-the-future:freeentry}{free-entry equilibrium}.
\item What distinguishes Sutton's two cases?
\end{enumerate}

\medskip\noindent\textbf{Part II --- The record.}
\begin{enumerate}[resume]
\item What did the SEC's 2022 release find about wholesaling, and what happened to the proposal?
\item State the 2023 merger guidelines' thresholds.
\item Describe the consolidation chain of \cref{fig:fm:competition-and-the-future:chain}.
\item What do the \omterm{def:fm:competition-and-the-future:token}{tokenised fund}'s prospectus, the 24X order and NSCC's announcement establish?
\end{enumerate}

\medskip\noindent\textbf{Part III --- The model.}
\begin{enumerate}[resume]
\item State and prove \cref{prop:fm:competition-and-the-future:bounds}; give the bounds for the wholesalers.
\item Derive the free-entry number of firms in the symmetric Cournot model.
\item Give the number of firms and the index at \$25, 50 and 100 million.
\item Give the planner's numbers and explain the gap.
\item Give the merger's effect on price, index and profits, and the synergy that holds the price.
\item Why does the model's index fall short of the public bounds?
\end{enumerate}

\medskip\noindent\textbf{Part IV --- The strategy.}
\begin{enumerate}[resume]
\item Which of the firm's fixed costs do cheaper tools lower, and which do round-the-clock markets raise?
\item In which of the firm's businesses is the fixed cost chosen (a spending race)?
\item How would you tell, from the firm's own numbers, whether its moat is shrinking?
\item What would you watch in the record over the next two years?
\item State the \emph{named result}: the free-entry number of firms and the index at the book's parameters, and how both move when the fixed cost
halves.
\item In two sentences, write the recommendation.
\end{enumerate}
\end{problem}

\section{Interview questions}

\begin{interviewq}[$\star$ \iqroles{researcher}]\label{iq:fm:competition-and-the-future:1}
What is the \omterm{def:fm:competition-and-the-future:hhi}{Herfindahl--Hirschman index} of a business with four equal firms, and what does a merger of two of them do to it?
\end{interviewq}

\begin{interviewq}[$\star$ \iqroles{trader, researcher}]\label{iq:fm:competition-and-the-future:2}
What is your firm's moat, if any?
\end{interviewq}

\begin{interviewq}[$\star\star$ \iqroles{researcher}]\label{iq:fm:competition-and-the-future:3}
In a Cournot market with linear demand and a fixed cost, how many firms enter, and how does the number scale with the market's size?
\end{interviewq}

\begin{interviewq}[$\star\star$ \iqroles{developer}]\label{iq:fm:competition-and-the-future:4}
What changes in a trading system when the market it trades runs 23 hours a day?
\end{interviewq}

\begin{interviewq}[$\star\star$ \iqroles{risk}]\label{iq:fm:competition-and-the-future:5}
A counterparty offers tokenised money-market fund shares as collateral on a Sunday. What do you check?
\end{interviewq}

\begin{interviewq}[$\star\star\star$ \iqroles{researcher}]\label{iq:fm:competition-and-the-future:6}
You know only that six firms serve a business and that the two largest hold two thirds. Bound its concentration.
\end{interviewq}
